\subsubsection*{Statement B.2p (Refuted auxiliary lemma): a graph with an elliptic target}
\begin{proposition}
\label{B:B2-abelian-graph-counterexample}
There is a geometrically normal projective integral surface $M$
over an imperfect field $k$, an elliptic curve $E/k$, and a rational
map $M\dashrightarrow E$ for which no proper geometrically normal
modification of $M$ resolves the map.  In particular, the normalization
of a graph with an abelian target need not be geometrically normal.
\end{proposition}
\begin{proof}
We work in characteristic three and put $k=\mathbf F_3(a)$, where
$a$ is transcendental.  In $\mathbf P^2_k$ let
\[
 E:\quad F:=Y^2Z-X^3+XZ^2=0,\qquad Q:=Y^2+Z^2.
\]
The gradient of $F$ is
$(Z^2,-YZ,Y^2-XZ)$.  Its simultaneous vanishing on $F=0$
would first give $Z=0$, then $X=0$, then $Y=0$, which is
impossible in projective space.  Thus $E$ is a smooth plane cubic.
With origin $[0:1:0]$ it is an elliptic curve.  The divisor
$E\cap\{Q=0\}$ consists of six distinct geometric points.
Indeed none lies on $Z=0$, and in the chart $Z=1$ its equations
are
\[
                    y^2+1=0,\qquad x^3-x+1=0.
\]
Both polynomials are separable, and their degrees are two and
three.  The gradients of $F$ and $Q$ are independent at all
these points: the first has nonzero $X$-component $Z^2$,
whereas the second has zero $X$-component and nonzero
$Y$-component $-Y$.

Let $C\subset\mathbf P^3_k$, with coordinates $[X:Y:Z:W]$,
be the projective cone $F=0$.  Define $M\subset\mathbf P^4_k$
by
\[
          F=0,\qquad V^3=aW^3+WQ.
          \tag{B.2.7}
\]
Projection forgetting $V$ is a finite flat map $M\to C$ of
degree three: on each standard affine chart of $C$ the equation
for the appropriately scaled $V$ is monic of degree three.
The projection centre does not meet $M$.

We verify geometric normality directly.  Over an algebraic closure
$\bar k$, choose $b^3=a$ and write $V'=V-bW$.  The equations
become
\[
                         F=0,\qquad (V')^3=WQ.
\]
The coordinate ring of the affine cone over $E_{\bar k}$ is a
domain, and $WQ$ is not a cube in its rational function field:
its $W$-adic valuation is one.  Hence the latter equation is
irreducible over that function field.  Since it is monic, the
homogeneous coordinate ring of $M_{\bar k}$ is a domain.
The two displayed equations form a regular sequence, so
$M_{\bar k}$ is Cohen--Macaulay and in particular satisfies $S_2$.

At a point with $(X,Y,Z)\ne(0,0,0)$, the gradient of $F$ is
nonzero.  If $Q\ne0$, the $W$-component of the gradient of
$(V')^3-WQ$ is nonzero, so the two gradients are independent.
If $Q=0$, we have $Z\ne0$ and $Y\ne0$.  A proportionality
between the gradients would have proportionality factor zero,
by their $X$-components; it therefore forces $W=0$ and then
$V'=0$.  These are precisely the six points just described.
The remaining possible singular point is
$[0:0:0:1:0]$.  Thus the geometric singular locus is finite.
As $M_{\bar k}$ has dimension two, it is regular in codimension
one.  Serre's criterion now proves that $M$ is geometrically
normal, as asserted.

For completeness we describe the graph without invoking an
unspecified resolution.  The incidence variety
\[
 Y_0=\{(c,e)\in C\times E:
              (X,Y,Z)\text{ is proportional to the coordinates of }e\}
\]
is a ruled surface over $E$, isomorphic to
$\mathbf P_E(\mathcal O_E\oplus\mathcal O_E(1))$
with the quotient convention for projective bundles.  Its map to
$C$ is an isomorphism away from the cone vertex and contracts
one section $E_0$ to that vertex.  The displayed incidence
condition consists of the three $2\times2$ minors, so $Y_0$
is a closed subscheme of $C\times E$.
Put
\[
                         \Gamma=M\times_CY_0.
\]
It is finite flat over $Y_0$ and integral, by the same monic
generic-field argument used above.  It is a closed subscheme
of $M\times E$, and is proper birational over $M$.

Near the exceptional section choose a local trivialization of
$\mathcal O_E(1)$ and a transverse parameter $t$, so that the
cone coordinates are $t$ times the coordinates on $E$ and
$W=1$.  The local equation of $\Gamma$ is
\[
                         v^3=a+t^2q,
                         \tag{B.2.8}
\]
where $q$ represents the section $Q$ in this trivialization.
The equation $t=0$ cuts out the Cartier divisor
$D=E_l$, where $l=k(a^{1/3})$.  This is a regular scheme:
$E_l$ is smooth over the field $l$, including at its closed
points whose residue fields may be inseparable over $k$.
Since $t$ is a nonzerodivisor, regularity of the quotient by
$t$ implies regularity of $\Gamma$ along $D$.  Away from $D$,
the map $\Gamma\to M$ is an isomorphism.  Therefore $\Gamma$
is normal everywhere.  We do not claim that it is regular
at the six other singular points.

Because $\Gamma$ is integral and contains the graph over the
dense open on which the projection $M\dashrightarrow E$ is
defined, its closed incidence embedding identifies it with
the graph closure.  It is already normal, so it is also the
normalized graph.

After extension to $\bar k$, the local equation at the generic
point of the reduced exceptional divisor is
\[
            s^3=t^2q,\qquad q\in\bar k(E)^\times,
            \quad s=v-b.
\]
Its local ring at $(s,t)$ is a one-dimensional local domain
of embedding dimension two, hence is not a discrete valuation
ring and is not normal.  One can also see the defect explicitly:
$h=s^2/t$ is integral because $h^3=tq^2$, and
$t=h^3/q^2$, $s=h^2/q$.  The cusp ring has value semigroup
generated by two and three, whereas $h$ has value one, so
$h$ does not belong to it.

Finally, suppose a proper geometrically normal modification
$N\to M$ resolved the rational map to $E$.  The resulting
map $N\to M\times E$ factors through the graph $\Gamma$;
it is proper and birational.  At the generic point of the
$k$-divisor $D$, the scheme $\Gamma$ is regular of dimension
one.  A proper birational morphism between normal integral
schemes is an isomorphism over a neighbourhood of such a
codimension-one point.  Thus an open neighbourhood in $N$
would be isomorphic to an open neighbourhood of that point
in $\Gamma$.  After passage to $\bar k$ this neighbourhood
contains the nonnormal codimension-one ring just computed,
contrary to geometric normality of $N$.

The accompanying exact certificate checks the homogeneous
derivatives, the affine singular-locus ideal, the six distinct
points over $\mathbf F_{3^6}$, and the integral-element
identities.  The regularity and proper-birational arguments
above supply the geometric conclusions.
\end{proof}

\begin{remark}[What this counterexample resolves]
The ordinary normal graph always exists, and over a perfect
field its normality is geometric normality.  The extra hypothesis
that the target is an abelian variety does not provide this
conclusion over an imperfect field.  Thus the proposed abelian
graph construction cannot close B.2 in its original-field scope.
This proposition is not a counterexample to weak upper-cone
Picard stabilization; indeed B.2o proves that weak assertion for
the surface in (B.2.7).
\end{remark}
